\documentclass{article}

\PassOptionsToPackage{numbers, compress}{natbib}
\usepackage[preprint]{neurips_2026}

\usepackage[utf8]{inputenc}
\usepackage[T1]{fontenc}
\usepackage{hyperref}
\usepackage{url}
\usepackage{booktabs}
\usepackage{amsfonts}
\usepackage{amsmath}
\usepackage{nicefrac}
\usepackage{microtype}
\usepackage{xcolor}
\usepackage{graphicx}
\usepackage{subcaption}

\newcommand{\RQ}[1]{\textbf{RQ#1}}

\title{How Far Can 0.6B Parameters Go?\\Compute-Efficient Generative Question Answering\\with LoRA Fine-Tuning}

\author{%
  Author One \And Author Two \And Author Three \And Author Four \And Author Five \\
  EN.737 Large Language Models, Johns Hopkins University \\
  \texttt{\url{https://github.com/Syndr0/compute-efficient-genqa}}
}

\begin{document}

\maketitle

\begin{abstract}
%%ABSTRACT%%
\end{abstract}

\section{Introduction}
\label{sec:intro}

Small generative language models are attractive for question answering (QA) over a provided
passage: they run on a single commodity GPU, can be adapted cheaply, and produce free-form
answers without task-specific span-extraction heads. How much compute such adaptation needs,
how well the adapted model transfers to unseen domains, and how to serve it efficiently are
practical questions that are rarely answered together under one controlled setup.

We study these questions with \textbf{Qwen3-0.6B}~\citep{qwen3} adapted by low-rank adaptation
(LoRA)~\citep{hu2022lora} to context-grounded, single-turn QA on the MRQA benchmark
collection~\citep{fisch2019mrqa}. Every experiment shares one model, one prompt, one LoRA
configuration and one evaluation pipeline, and the training data is controlled by an exact
\emph{training-token budget}. We ask:
\begin{itemize}
  \item \RQ{1} (training efficiency): how does QA quality change with the LoRA training-token budget, relative to zero- and few-shot prompting?
  \item \RQ{2} (generalization): at a fixed budget, does multi-source fine-tuning (SQuAD + NewsQA + TriviaQA) generalize better to unseen domains (BioASQ, DuoRC) than SQuAD-only fine-tuning?
  \item \RQ{3} (inference efficiency): how does the inference batch size trade throughput against latency and GPU memory for the fine-tuned model?
\end{itemize}
Hypotheses and quantitative predictions were fixed before the main runs and are evaluated in
Appendix~\ref{app:hypotheses}.

%%CONTRIBUTIONS%%

\section{Related work}
\label{sec:related}

\paragraph{Parameter-efficient adaptation.} LoRA~\citep{hu2022lora} trains low-rank updates to
frozen weight matrices and is the default way to adapt LLMs on limited hardware; QLoRA~\citep{dettmers2023qlora}
further quantizes the frozen base. Prior work shows that a small amount of high-quality supervised
data can suffice to change a model's output behaviour~\citep{zhou2023lima}; we quantify how much
data is needed for extractive-style QA in terms of training tokens rather than examples.

\paragraph{Generalization in reading comprehension.} The MRQA shared task~\citep{fisch2019mrqa}
standardised six training and six held-out QA datasets to measure out-of-domain transfer. Training
on several sources was found to improve robustness of extractive models~\citep{talmor2019multiqa,fisch2019mrqa},
and multi-format training benefits generative QA models~\citep{khashabi2020unifiedqa,raffel2020t5}.
These studies typically compare mixtures of different total sizes; we hold the token budget fixed so
that the comparison isolates data diversity from data quantity.

\paragraph{Compute budgets and inference.} Scaling studies relate loss to parameters and training
tokens~\citep{kaplan2020scaling,hoffmann2022chinchilla}. For serving, batching amortizes the
weight reads of memory-bound autoregressive decoding~\citep{pope2023inference}, and systems such as
Orca~\citep{yu2022orca} and vLLM~\citep{kwon2023vllm} schedule requests dynamically. We measure the
basic static-batching trade-off that these systems build on.

\section{Experimental setup}
\label{sec:setup}

\paragraph{Data.} We use the MRQA 2019 versions~\citep{fisch2019mrqa} of SQuAD~\citep{rajpurkar2016squad},
NewsQA~\citep{trischler2017newsqa} and TriviaQA~\citep{joshi2017triviaqa} (web) for training and
in-domain (ID) evaluation, and BioASQ~\citep{tsatsaronis2015bioasq} and DuoRC
(ParaphraseRC)~\citep{saha2018duorc} as out-of-domain (OOD) test sets that are never used for
training or model selection. Training pools are capped at 50k questions per source. From each pool we
hold out whole passages (300 questions) as a development set used only for validation loss, so no
passage appears in both. ID test sets are drawn from the MRQA validation splits and OOD sets from
the MRQA test split; eval questions whose passage occurs verbatim (after normalisation) in a training
pool are removed before sub-sampling 1,000 questions per set (Table~\ref{tab:data}). MRQA markup
tokens are converted to plain text.

\paragraph{Task format and model.} Each input is one chat turn: a system instruction asking for the
shortest answer span copied from the passage, followed by \texttt{Passage: \dots\ Question: \dots}.
Qwen3's thinking mode is disabled and passages are truncated to 768 tokens (at test time the head
of the passage is kept; for training the window is moved to contain the answer). Few-shot prompting
prepends three fixed short SQuAD training examples as earlier chat turns~\citep{brown2020gpt3}.
Answers are generated greedily (at most 32 tokens) and the first line is kept.

\paragraph{LoRA fine-tuning under a token budget.} LoRA adapters (rank 16, $\alpha{=}32$, dropout
0.05) are attached to all attention and MLP projections, giving 10.1M trainable parameters (1.7\% of
the model). We train with AdamW, learning rate $2{\times}10^{-4}$, 5\% warm-up and cosine decay,
16 sequences per optimizer step and the loss on answer tokens only. A budget $B$ means that examples
are drawn without replacement---with an equal token share per source---until their prompt plus answer
tokens sum to $B$; the model makes a single pass with a schedule fitted to that length. Training cost
is reported as wall-clock GPU time and estimated FLOPs $\approx 4NT$ for $N$ model parameters and $T$
tokens (forward plus activation gradients; the frozen base needs no weight gradients).
\RQ{1} trains on the three-source mixture with $B \in \{0.25, 1, 4, 16\}$M tokens.
\RQ{2} compares SQuAD-only against the three-source mixture at $B{=}4$M with three seeds (data
sample and adapter initialisation), sharing all other settings.

\paragraph{Evaluation and statistics.} We report Exact Match (EM) and token-level F1 with the official
SQuAD/MRQA normalisation, maximised over gold answers. Confidence intervals are 95\% percentile
bootstrap intervals over questions~\citep{efron1993bootstrap}; system comparisons use a paired
bootstrap over the same questions~\citep{koehn2004bootstrap} (10,000 resamples), with per-question
scores averaged over seeds~\citep{dodge2020finetuning}. Errors are categorised automatically
(\S\ref{sec:rq2}).

\paragraph{Inference benchmark.} \RQ{3} decodes the same 512 SQuAD prompts with the merged
4M-token adapter in bf16 using HuggingFace \texttt{generate}~\citep{wolf2020transformers}:
batch sizes 1--256 with length-sorted, left-padded batches and the KV cache on, plus ablations with
unsorted batches, the KV cache disabled, and fp32/fp16 weights. Each configuration has one
warm-up batch and three timed repeats; quality is recorded to confirm that batching does not change
the answers.

\paragraph{Hardware and software.} All final runs used one NVIDIA RTX PRO 6000 Blackwell
Server Edition GPU (96~GB; Google Colab ``G4'') with PyTorch, Transformers~\citep{wolf2020transformers}
and PEFT~\citep{peft}. The full pipeline, configurations, per-question predictions and logs are released
with the paper.

\section{Results}
\label{sec:results}

%%RESULTS%%

\section{Discussion and limitations}
\label{sec:discussion}

%%DISCUSSION%%

\section{Conclusion}
\label{sec:conclusion}

%%CONCLUSION%%

\paragraph{Reproducibility.} All numbers in this paper are produced by
\texttt{python -m cegqa.cli all --config configs/default.yaml --config configs/a100.yaml}
from the released repository, which also writes the tables and figures used here. The notebooks in
\texttt{notebooks/} run the same pipeline on Google Colab.

\paragraph{Author contributions.} P1 built the data and evaluation pipeline (MRQA processing, passage-level
splits, overlap checks, EM/F1 and prediction logging). P2 built the prompting baselines and the LoRA training
pipeline with token-budget sampling and cost accounting (\RQ{1}). P3 designed and analysed the generalization
study, including the paired bootstrap and error analysis (\RQ{2}). P4 built the inference benchmark (\RQ{3}).
P5 led the analysis, figures and writing. All authors interpreted the results and edited the paper.

\bibliographystyle{plainnat}
\bibliography{references}

\appendix

%%APPENDIX%%

\end{document}
